\documentclass[12pt,a4paper]{ctexart}

\usepackage[margin=2.2cm]{geometry}
\usepackage{amsmath,amssymb}
\usepackage{array,booktabs,tabularx}
\usepackage{enumitem}
\usepackage[dvipsnames]{xcolor}
\usepackage{hyperref}

\hypersetup{
  colorlinks=true,
  linkcolor=blue!60!black,
  urlcolor=blue!60!black
}

\title{3.3.1 极简版 Self-Attention 讲义\\\large 先问题，再直觉，再公式，再代码}
\author{}
\date{\today}

\begin{document}

\maketitle
\tableofcontents
\newpage

\section{这一节到底在解决什么}

这一节不是在讲完整的 Transformer，而是在讲它里面最核心的一根骨架：

\begin{quote}
当我们更新一个 token 的表示时，怎样让它有选择地吸收其他 token 的信息？
\end{quote}

作者先故意讲一个\textbf{没有可训练参数}的版本，不是因为真实模型就这么做，而是因为这样最容易把 attention 的本质看清楚。

这一节真正想让你先看懂的，不是复杂公式，而是下面这条链：

\[
\text{相关性打分} \rightarrow \text{softmax 归一化} \rightarrow \text{加权求和}
\]

\section{新的讲解原则}

从现在开始，每个知识点都按下面四层来讲：

\begin{enumerate}[leftmargin=2em]
  \item \textbf{背景问题}：它要解决什么现实困难；
  \item \textbf{朴素直觉}：先用一句人话说清楚它在做什么；
  \item \textbf{数学表达}：再给公式，而且每个符号都翻译成人话；
  \item \textbf{代码对应}：最后再落回 notebook 里的具体代码。
\end{enumerate}

这一份关于 self-attention 的讲义，也严格按这个顺序展开。

\section{新手最该先问的 5 个问题}

\begin{enumerate}[leftmargin=2em]
  \item 什么叫``更新一个 token 的表示''？
  \item 为什么这里会冒出一个新的向量 \(z^{(i)}\)？
  \item 为什么先只看第二个词 \(x^{(2)}\)？
  \item 为什么相关性可以用点积来表示？
  \item 为什么最后是加权求和，而不是只挑最相关的那个词？
\end{enumerate}

\section{QA 讲解}

\subsection{Q1：什么叫更新一个 token 的表示？}

答：原始输入里的每个 token 已经先有一个 embedding，记成 \(x^{(i)}\)。
但这个 embedding 只说明``它自己是谁''，还没有显式融合上下文。

attention 要做的事是：围绕某个 token，把整句话里和它相关的信息重新汇总一遍，得到一个新的表示 \(z^{(i)}\)。

所以可以先把 \(z^{(i)}\) 理解成：

\begin{quote}
原来的 token 表示，加上它从其他 token 那里收集来的上下文信息。
\end{quote}

\subsection{Q2：为什么这里会冒出一个新的上下文向量？}

答：因为模型真正需要的，不只是``这个词本身是什么''，还需要``这个词在当前句子里意味着什么''。

例如同一个词放在不同句子里，依赖的上下文会不同。
所以 self-attention 的输出不是把 \(x^{(i)}\) 原样抄一遍，而是生成一个\textbf{带上下文的新表示} \(z^{(i)}\)。

\subsection{Q3：为什么先只看第二个词？}

答：因为如果一上来就把所有 \(z^{(1)}, z^{(2)}, \dots, z^{(T)}\) 全部同时写出来，新手很容易被下标淹没。

作者先固定一个具体例子：第二个词 \texttt{journey}，记作 \(x^{(2)}\)。
这样问题就收缩成一句话：

\begin{quote}
怎样根据 \(x^{(2)}\)，从整句话里挑出对它最有用的信息？
\end{quote}

等你看懂一个 token 的流程后，再推广到全部 token 就很自然了。

\subsection{Q4：为什么相关性可以用点积来表示？}

答：在这个极简版里，每个 token 都是一个向量。
两个向量做点积时，如果它们的方向更一致、各维度上更同向，结果通常更大。

所以这里把点积当成一个粗略的``相似性分数''：

\[
\omega_{2i} = x^{(i)} q^{(2)T}
\]

其中 \(q^{(2)} = x^{(2)}\)。

这句话的人话翻译是：

\begin{quote}
把第二个词当作当前关注对象，看句子里每个词和它有多对齐。
\end{quote}

注意：这里的 \(\omega\) 只是原始分数，还不是最终权重。

\subsection{Q5：为什么最后是加权求和，而不是只挑最相关的那个词？}

答：因为一个 token 往往不是只依赖一个位置，而是可能同时依赖多个位置。

如果只取最大值，相当于说``我只看最相关的一个词''，这会丢掉其他有用信息。
attention 的做法更柔和：

\begin{enumerate}[leftmargin=2em]
  \item 先给所有位置一个相关性分数；
  \item 再把分数变成总和为 1 的权重；
  \item 最后把所有输入向量按权重混合起来。
\end{enumerate}

数学上写成：

\[
\alpha_{2i} = \frac{e^{\omega_{2i}}}{\sum_j e^{\omega_{2j}}}
\]

\[
z^{(2)} = \sum_i \alpha_{2i} x^{(i)}
\]

这就是这一节最核心的骨架。

\section{把 notebook 里的例子完整走一遍}

\subsection{输入是什么}

notebook 用一句很短的 toy sentence：

\begin{center}
\texttt{Your journey starts with one step}
\end{center}

每个词先被表示成一个 3 维向量：

\begin{center}
\begin{tabular}{c c c}
\toprule
token & 记号 & 向量 \\
\midrule
Your & \(x^{(1)}\) & \([0.43, 0.15, 0.89]\) \\
journey & \(x^{(2)}\) & \([0.55, 0.87, 0.66]\) \\
starts & \(x^{(3)}\) & \([0.57, 0.85, 0.64]\) \\
with & \(x^{(4)}\) & \([0.22, 0.58, 0.33]\) \\
one & \(x^{(5)}\) & \([0.77, 0.25, 0.10]\) \\
step & \(x^{(6)}\) & \([0.05, 0.80, 0.55]\) \\
\bottomrule
\end{tabular}
\end{center}

这里 3 维只是为了演示方便，真实模型里的维度会大很多。

\subsection{第 1 步：算 attention scores}

这一节把第二个词当作查询：

\[
q^{(2)} = x^{(2)} = [0.55, 0.87, 0.66]
\]

然后和所有输入向量做点积：

\[
\omega_{21} = x^{(1)} \cdot q^{(2)} = 0.9544
\]
\[
\omega_{22} = x^{(2)} \cdot q^{(2)} = 1.4950
\]
\[
\omega_{23} = x^{(3)} \cdot q^{(2)} = 1.4754
\]
\[
\omega_{24} = 0.8434,\quad
\omega_{25} = 0.7070,\quad
\omega_{26} = 1.0865
\]

把它们合起来就是：

\[
[0.9544,\ 1.4950,\ 1.4754,\ 0.8434,\ 0.7070,\ 1.0865]
\]

直觉上可以这样读：

\begin{quote}
对于当前的 \texttt{journey}，\texttt{journey} 自己和 \texttt{starts} 的分数比较高，说明它们和当前查询更对齐。
\end{quote}

\subsection{第 2 步：把 scores 变成真正的权重}

这些分数还不能直接拿来加权，因为它们的和不是 1，而且也不方便解释成``占比''。

所以作者接着做 softmax，得到 attention weights：

\[
[\alpha_{21}, \alpha_{22}, \alpha_{23}, \alpha_{24}, \alpha_{25}, \alpha_{26}]
= [0.1385,\ 0.2379,\ 0.2333,\ 0.1240,\ 0.1082,\ 0.1581]
\]

现在这些数可以直接理解成：

\begin{quote}
当我们更新 \(x^{(2)}\) 时，句子里每个词各自应该贡献多少比例。
\end{quote}

它们的和等于 1，这正是 softmax 的作用。

\subsection{第 3 步：做加权求和，得到上下文向量}

最后把所有输入向量按这些权重做加权平均：

\[
z^{(2)} = \sum_i \alpha_{2i} x^{(i)}
\]

在 notebook 的例子里，算出来是：

\[
z^{(2)} = [0.4419,\ 0.6515,\ 0.5683]
\]

这个 \(z^{(2)}\) 就是第二个词的新表示。
它不再只是原始的 \texttt{journey} embedding，而是已经混入了整句话里其他位置的信息。

\section{这一节和真实 Transformer 还差什么}

最关键的一点是：这一版\textbf{没有可训练参数}。

这里直接用了：

\begin{enumerate}[leftmargin=2em]
  \item \(q^{(2)} = x^{(2)}\)
  \item 被比较的对象也是原始输入 \(x^{(i)}\)
  \item 最后加权的还是原始输入 \(x^{(i)}\)
\end{enumerate}

而真正的 Transformer 不会直接拿原始输入互相比，它会先学出三套变换：

\[
Q = XW_Q,\quad K = XW_K,\quad V = XW_V
\]

也就是：

\begin{enumerate}[leftmargin=2em]
  \item 用 \(Q\) 表示``我现在想找什么''；
  \item 用 \(K\) 表示``我这里有什么可供匹配''；
  \item 用 \(V\) 表示``如果你来取我，你最终取走什么信息''。
\end{enumerate}

所以这一节不是完整答案，而是把完整答案拆开前的准备动作。

\section{代码里每一步对应什么}

文件位置是：

\begin{center}
\texttt{ch03/01\_main-chapter-code/ch03.ipynb}
\end{center}

这一小节可以按下面的顺序看代码：

\begin{enumerate}[leftmargin=2em]
  \item \texttt{inputs = torch.tensor(...)}：定义 6 个 token 的 toy embedding。
  \item \texttt{query = inputs[1]}：固定第二个词作为当前查询。
  \item \texttt{torch.dot(x\_i, query)}：计算每个词和查询的点积分数。
  \item \texttt{softmax\_naive} 或 \texttt{torch.softmax}：把分数变成和为 1 的权重。
  \item \texttt{context\_vec\_2 += attn\_weights\_2[i] * x\_i}：做加权求和，得到 \(z^{(2)}\)。
\end{enumerate}

如果你是第一次看 notebook，真正要盯住的是变量流向：

\[
inputs \rightarrow query \rightarrow attn\_scores \rightarrow attn\_weights \rightarrow context\_vec
\]

\section{这一节真正该记住什么}

这一节最重要的不是记住具体数字，而是记住下面这句话：

\begin{quote}
self-attention 的最小骨架就是：选定一个 token，计算它和所有 token 的相关性，把相关性变成权重，再用这些权重对整句信息做加权汇总。
\end{quote}

如果再压缩成更短的一句话，就是：

\begin{quote}
相关性打分 $\rightarrow$ 归一化 $\rightarrow$ 加权求和。
\end{quote}

\end{document}
